\section{Limits for Each Protocol}
\label{sec:limits}

This section shows two limits. They apply to each protocol in the model of
\cref{sec:model}, and not only to \sys{}. The two limits come from one lemma.

\subsection{The Substitution Lemma}

\begin{lem}[Substitution]
\label{lem:subst}
Let $\rho$ be a run of a protocol in which less than $f$ officers are hostile. Let $x$ be
an officer who is loyal in $\rho$, and let $d$ be a device. Let $t_0$ be a time such that,
from $t_0$, $x$ and $d$ receive only messages that $x$ or $d$ sent. Then the protocol has
a run $\rho'$ with these properties.
\begin{enumerate}
  \item Until $t_0$, the run $\rho'$ is the same as $\rho$.
  \item At $t_0$, the attacker gets the key and the data of $x$. It gets no other key that
  it does not have in $\rho$.
  \item From $t_0$, each party other than $x$ and $d$ does the steps that it selects from
  the steps that the protocol lets it do.
  \item At all times, the view of $d$ in $\rho'$ is the same as the view of $d$ in $\rho$.
\end{enumerate}
\end{lem}

\begin{proof}
Let each clock in $\rho'$ be the same function of real time as in $\rho$. From $t_0$ in
$\rho'$, the attacker does three steps. First, it stops each message that a party sends to
$x$ or to $d$. Second, it makes a copy of $x$ from the data and the key of
$x$. It gives the copy each message of $d$ that $x$ received in $\rho$, at the same time as
in $\rho$. Third, it gives $d$ each message that $d$ received in $\rho$, at the same time
as in $\rho$.

We show by induction on the sequence of steps that the attacker can make these messages.
Until $t_0$, the two runs are the same. Before a step after $t_0$, let $d$ and the copy
have the views that $d$ and $x$ have in $\rho$. A device is deterministic, and the attacker
cannot change its clock. Thus $d$ sends the same message as in $\rho$, at the same time.
The copy has the data and the key of $x$, and the attacker selects for it the decisions
that $x$ made in $\rho$.

Thus the copy makes each message that $x$ sent in $\rho$. In $\rho$, $x$ and $d$ receive
no message from the other parties after $t_0$. Thus no message of the other parties is
necessary for the attacker. The steps of the other parties in $\rho'$ change no message
that $d$ receives. The view of $d$ is the same in the two runs at all times.
\end{proof}

The lemma uses one fact about the device: messages are its only input. A person at the
device is an input that the attacker does not control. Only such an input can make
\cref{lem:subst} not applicable.

\subsection{The Limit for the Commander}

\begin{thm}[Tenure limit]
\label{thm:limit}
Let a protocol have outage tolerance $P$ at a device $d$. Then the protocol has hostile
tenure $P$ at $d$. Also, in the run that shows the hostile tenure, $d$ accepts the same
orders at the same times as in the run with no attacker.
\end{thm}

\begin{proof}
Let $\rho$ be the run of \cref{def:tolerance}, with the commander $c$ and the times $t_0$
and $t_1$. No officer is hostile in $\rho$, and $f \ge 1$, thus \cref{lem:subst} applies
with $x = c$. In the third property, let each other
officer do the steps that remove $c$. The run $\rho'$ has the properties of
\cref{def:tenure}. The view of $d$ is the same as in $\rho$. Thus $d$ accepts the same
orders at the same times, and one of them is at $t_1 \ge t_0 + P$.
\end{proof}

The theorem tells that the hostile tenure is not less than the outage tolerance. Four
points make the result more accurate.

\paragraph{The outage is sufficient for the attacker}
In the proof, the attacker stops the messages to $c$ and to $d$. In the run with no
attacker, a failure stopped the same messages. Thus control of all links is not
necessary for the attacker. The key of the commander and the same outage are sufficient.

\paragraph{Limits on the number of orders}
The proof does not use the fact that $P$ is a time. Let a device accept $N$ orders of an
isolated loyal commander. Then it accepts the same $N$ orders of a hostile commander. Thus
a limit on the number of orders has the same property as a limit on time.

\paragraph{Checks of the contents of an order}
The device accepts the same orders in the two runs. Thus a check of the contents cannot
identify the run with the attacker. Such a check can reject only the orders that a loyal
commander does not give. A loyal operator and an attacker can open the same breaker
(\cref{sec:background}).

\paragraph{A person at the device}
The limit applies when messages are the only input of the device. If a person at the
device can stop remote orders, the hostile tenure stops when that person does this. Rule
D3 of \sys{} gives this input (\cref{sec:protocol}).

\subsection{The Limit for the Deputy}

\begin{thm}[Timeout takeover]
\label{thm:deputy}
Let a protocol have deputy period $F$ at a device $d$. Then the protocol has deputy
exposure $F$ at $d$.
\end{thm}

\begin{proof}
Let $\rho$ be the run of \cref{def:deputy}, with the deputy $b$ and the times $t_0$, $t_1$
and $t_2$. All officers are loyal in $\rho$, thus \cref{lem:subst} applies with $x = b$.
In the third property, let the commander
and the other officers do the steps of usual operation. Let the attacker stop no message
between them.

The run $\rho'$ has the properties of \cref{def:exposure}. The view
of $d$ is the same as in $\rho$. Thus $d$ accepts the orders of $b$ at $t_1$ and at
$t_2 \ge t_1 + F$.
\end{proof}

The result has a second part. Let the protocol let a commander continue when all officers
other than one are loyal and connected. Then, in the run of the proof, a different device
accepts orders of the commander while $d$ accepts orders of the deputy. Two officers give
orders at the same time, and the protocol does not have unity of command.

\Cref{def:deputy} includes each rule that gives command to a deputy when messages stop.
Examples are a timeout with a specified length, and the end of a lease. The theorem thus
tells that each such rule gives its period to an attacker who has the key of the deputy.

\subsection{Results for a Design}
\label{sec:design}

The two theorems give four rules for a design.
\begin{enumerate}
  \item \emph{An end is necessary.} If a device obeys an isolated commander without a time
  limit, a hostile commander keeps the device without a time limit.
  \item \emph{Zero has a cost.} If the hostile tenure must be zero, the outage tolerance is
  zero. A design with a quorum for each order has this property.
  \item \emph{One number for each device.} The designer selects the time $H$ for which a
  device obeys an isolated commander. This time is also the minimum of the hostile tenure.
  The best protocol makes $H$ also the maximum, with hostile officers, with stopped
  messages, and with no message to the device.
  \item \emph{A deputy period is an exposure.} The designer selects the period $F$ that a
  deputy gets without a quorum. A hostile deputy gets the same period. The safe selection
  is $F = 0$.
\end{enumerate}

The proofs are short. The argument is the argument of the CAP
result~\cite{gilbert2002cap}. A party that has the same view in two runs does the same
steps in the two runs. The theorems give the accurate relation for command. The minimum is the outage tolerance, and it applies to each protocol.
